% problem_id: S001-theorem-dold-kan
% source_id: S001 (Stacks Project)
% source_file: simplicial.tex
% statement_locator: simplicial.tex lines 4044--4055
% proof_locator: simplicial.tex lines 4057--4273
% env: theorem
% status: text-extracted-verbatim (difficulty judgement pending, written by hand)
%
% The statement and proof below are the author's own text, extracted verbatim.
% Nothing is paraphrased, shortened, or supplemented.

\begin{theorem}
\label{theorem-dold-kan}
Let $\mathcal{A}$ be an abelian category.
The functor $N$ induces an equivalence of
categories
$$
N :
\text{Simp}(\mathcal{A})
\longrightarrow
\text{Ch}_{\geq 0}(\mathcal{A})
$$
\end{theorem}

\begin{proof}
We will describe a functor in the reverse direction
inspired by the construction of Lemma \ref{lemma-extension}
(except that we throw in a sign to get the boundaries
right). Let $A_\bullet$ be a chain complex with boundary maps
$d_{A, n} : A_n \to A_{n - 1}$. For each $n \geq 0$ denote
$$
I_n =
\Big\{
\alpha : [n] \to \{0, 1, 2, \ldots\}
\mid
\Im(\alpha) = [k]\text{ for some }k
\Big\}.
$$
For $\alpha \in I_n$ we denote $k(\alpha)$ the unique
integer such that $\Im(\alpha) = [k]$.
We define a simplicial object $S(A_\bullet)$ as follows:
\begin{enumerate}
\item $S(A_\bullet)_n = \bigoplus_{\alpha \in I_n} A_{k(\alpha)}$, which
we will write as
$\bigoplus_{\alpha \in I_n} A_{k(\alpha)} \cdot \alpha$
to suggest thinking of ``$\alpha$'' as a basis vector for the
summand corresponding to it,
\item given $\varphi : [m] \to [n]$ we define
$S(A_\bullet)(\varphi)$ by its restriction to
the direct summand $A_{k(\alpha)} \cdot \alpha$
of $S(A_\bullet)_n$ as follows
\begin{enumerate}
\item $\alpha \circ \varphi \not \in I_m$ then we set it equal to zero,
\item $\alpha \circ \varphi \in I_m$ but $k(\alpha \circ \varphi)$
not equal to either $k(\alpha)$ or $k(\alpha) - 1$ then we set it
equal to zero as well,
\item if $\alpha \circ \varphi \in I_m$
and $k(\alpha \circ \varphi) = k(\alpha)$ then we use
the identity map to the summand
$A_{k(\alpha \circ \varphi)} \cdot (\alpha \circ \varphi)$
of $S(A_\bullet)_m$, and
\item if $\alpha \circ \varphi \in I_m$
and $k(\alpha \circ \varphi) = k(\alpha) - 1$
then we use $(-1)^{k(\alpha)} d_{A, k(\alpha)}$ to the summand
$A_{k(\alpha \circ \varphi)}\cdot (\alpha \circ \varphi)$
of $S(A_\bullet)_m$.
\end{enumerate}
\end{enumerate}
Let us show that $S(A_\bullet)$ is a simplicial object of $\mathcal{A}$.
To do this, assume we have maps $\varphi : [m] \to [n]$ and
$\psi : [n] \to [p]$. We will show that
$S(A_\bullet)(\varphi) \circ S(A_\bullet)(\psi) =
S(A_\bullet)(\psi \circ \varphi)$. Choose $\beta \in I_p$ and set
$\alpha = \beta \circ \psi$ and $\gamma = \alpha \circ \varphi$
viewed as maps $\alpha : [n] \to \{0, 1, 2, \ldots\}$
and $\gamma : [m] \to \{0, 1, 2, \ldots\}$. Picture
$$
\xymatrix{
[m] \ar[r]_\varphi \ar[d]_\gamma &
[n] \ar[r]_\psi \ar[d]_\alpha &
[p] \ar[d]^\beta \\
\Im(\gamma) \ar[r] &
\Im(\alpha) \ar[r] &
[k(\beta)]
}
$$
We will show that the restriction of the maps
$S(A_\bullet)(\varphi) \circ S(A_\bullet)(\psi)$ and
$S(A_\bullet)(\psi \circ \varphi)$.
to the summand $A_{k(\beta)} \cdot \beta$ agree.
There are several cases to consider
\begin{enumerate}
\item Say $\alpha \not \in I_n$ so the restriction of
$S(A_\bullet)(\psi)$ to $A_{k(\beta)} \cdot \beta$ is zero.
Then either $\gamma \not \in I_m$ or we have
$[k(\gamma)] = \Im(\gamma) \subset \Im(\alpha) \subset [k(\beta)]$
and the subset $\Im(\alpha)$ of $[k(\beta)]$ has a gap so
$k(\gamma) < k(\beta) - 1$. In both cases we see that
the restriction of
$S(A_\bullet)(\psi \circ \varphi)$ to $A_{k(\beta)} \cdot \beta$
is zero as well.
\item Say $\alpha \in I_n$ and $k(\alpha) < k(\beta) - 1$ so the restriction of
$S(A_\bullet)(\psi)$ to $A_{k(\beta)} \cdot \beta$ is zero.
Then either $\gamma \not \in I_m$ or we have
$[k(\gamma)] \subset [k(\alpha)] \subset [k(\beta)]$
and it follows that $k(\gamma) < k(\beta) - 1$. In both cases we see that
the restriction of
$S(A_\bullet)(\psi \circ \varphi)$ to $A_{k(\beta)} \cdot \beta$
is zero as well.
\item Say $\alpha \in I_n$ and $k(\alpha) = k(\beta)$ so the restriction of
$S(A_\bullet)(\psi)$ to $A_{k(\beta)} \cdot \beta$ is
the identity map from $A_{k(\beta)} \cdot \beta$ to
$A_{k(\alpha)} \cdot \alpha$. In this case because $\Im(\alpha) =[k(\beta)]$
the rule describing the restriction of $S(A_\bullet)(\psi \circ \varphi)$
to the summand $A_{k(\beta)} \cdot \beta$ is exactly the same
as the rule describing the restriction of $S(A_\bullet)(\varphi)$
to the summand $A_{k(\alpha)} \cdot \alpha$ and hence agreement holds.
\item Say $\alpha \in I_n$ and $k(\alpha) = k(\beta) - 1$
so the restriction of $S(A_\bullet)(\psi)$ to $A_{k(\beta)} \cdot \beta$
is given by $(-1)^{k(\beta)}d_{A, k(\beta)}$ to $A_{k(\alpha)} \cdot \alpha$.
Subcases
\begin{enumerate}
\item If $\gamma \not \in I_m$, then both the restriction of
$S(A_\bullet)(\psi \circ \varphi)$ to the summand $A_{k(\beta)} \cdot \beta$
and the restriction of $S(A_\bullet)(\varphi)$ to the summand
$A_{k(\alpha)} \cdot \alpha$ are zero and we get agreement.
\item If $\gamma \in I_m$ but $k(\gamma) < k(\alpha) - 1$, then
again both restrictions are zero and we get agreement.
\item If $\gamma \in I_m$ and $k(\gamma) = k(\alpha)$ then
$\Im(\gamma) = \Im(\alpha)$. In this case
the restriction of $S(A_\bullet)(\psi \circ \varphi)$
to the summand $A_{k(\beta)} \cdot \beta$ is given by
$(-1)^{k(\beta)}d_{A, k(\beta)}$ to $A_{k(\gamma)} \cdot \gamma$
and the restriction of $S(A_\bullet)(\varphi)$ to the summand
$A_{k(\alpha)} \cdot \alpha$ is the identity map
$A_{k(\alpha)} \cdot \alpha \to A_{k(\gamma)} \cdot \gamma$.
Hence agreement holds.
\item Finally, if $\gamma \in I_m$ and $k(\gamma) = k(\alpha) - 1$ then
the restriction of $S(A_\bullet)(\varphi)$ to the summand
$A_{k(\alpha)} \cdot \alpha$ is given by $(-1)^{k(\alpha)} d_{A, k(\alpha)}$
as a map
$A_{k(\alpha)} \cdot \alpha \to A_{k(\beta)} \cdot \beta$.
Since $A_\bullet$ is a complex we see that the composition
$A_{k(\beta)} \cdot \beta \to
A_{k(\alpha)} \cdot \alpha \to A_{k(\gamma)} \cdot \gamma$ is zero
which matches what we get for the restriction of
$S(A_\bullet)(\psi \circ \varphi)$
to the summand $A_{k(\beta)} \cdot \beta$
because $k(\gamma) = k(\beta) - 2 < k(\beta) - 1$.
\end{enumerate}
\end{enumerate}
Thus $S(A_\bullet)$ is a simplicial object of $\mathcal{A}$.

\medskip\noindent
Let us construct an isomorphism $A_\bullet \to N(S(A_\bullet))$
functorial in $A_\bullet$. Recall that
$$
S(A_\bullet) =  N(S(A_\bullet)) \oplus D(S(A_\bullet))
$$
as chain complexes by Lemma \ref{lemma-decompose-associated-complexes}.
On the other hand it follows from Remark \ref{remark-degenerate-subcomplex}
and the construction of $S(A_\bullet)$ that
$$
D(S(A_\bullet))_n =
\bigoplus\nolimits_{\alpha \in I_n,\ k(\alpha) < n} A_{k(\alpha)} \cdot \alpha
\subset
\bigoplus\nolimits_{\alpha \in I_n} A_{k(\alpha)} \cdot \alpha
$$
However, if $\alpha \in I_n$ then we have
$k(\alpha) \geq n \Leftrightarrow \alpha = \text{id}_{[n]} : [n] \to [n]$.
Thus the summand $A_n \cdot \text{id}_{[n]}$ of $S(A_\bullet)_n$
is a complement to the summand $D(S(A_\bullet))_n$.
All the maps $d^n_i : S(A_\bullet)_n \to S(A_\bullet)_n$
restrict to zero on the summand $A_n \cdot \text{id}_{[n]}$
except for $d^n_n$ which produces $(-1)^n d_{A, n}$
from $A_n \cdot \text{id}_{[n]}$ to $A_{n - 1} \cdot \text{id}_{[n - 1]}$.
We conclude that $A_n \cdot \text{id}_{[n]}$ must be
equal to the summand $N(S(A_\bullet))_n$ and moreover
the restriction of the differential
$d_n = \sum (-1)^id^n_i : S(A_\bullet)_n \to S(A_\bullet)_{n - 1}$
to the summand $A_n \cdot \text{id}_{[n]}$ gives what we want!

\medskip\noindent
Finally, we have to show that $S \circ N$ is isomorphic to the
identity functor. Let $U$ be a simplicial object of $\mathcal{A}$.
Then we can define an obvious map
$$
S(N(U))_n = \bigoplus\nolimits_{\alpha \in I_n} N(U)_{k(\alpha)} \cdot \alpha
\longrightarrow
U_n
$$
by using $U(\alpha) : N(U)_{k(\alpha)} \to U_n$ on the summand corresponding
to $\alpha$. By Definition \ref{definition-split} this is an isomorphism.
To finish the proof we have to show that this is compatible with
the maps in the simplicial objects. Thus let $\varphi : [m] \to [n]$
and let $\alpha \in I_n$. Set $\beta = \alpha \circ \varphi$.
Picture
$$
\xymatrix{
[m] \ar[r]_\varphi \ar[d]_\beta &
[n] \ar[d]_\alpha \\
\Im(\beta) \ar[r] &
[k(\alpha)]
}
$$
There are several cases to consider
\begin{enumerate}
\item Say $\beta \not \in I_m$. Then there exists an index
$0 \leq j < k(\alpha)$ with $j \not \in \Im(\alpha \circ \varphi)$
and hence we can choose a factorization
$\alpha \circ \varphi = \delta^{k(\alpha)}_j \circ \psi$
for some $\psi : [m] \to [k(\alpha) - 1]$.
It follows that $U(\varphi)$ is zero on the image of the summand
$N(U)_{k(\alpha)} \cdot \alpha$ because $U(\varphi) \circ U(\alpha) =
U(\alpha \circ \varphi) = U(\psi) \circ d^{k(\alpha)}_j$
is zero on $N(U)_{k(\alpha)}$ by construction of $N$.
This matches our rule for $S(N(U))$ given above.
\item Say $\beta \in I_m$ and $k(\beta) < k(\alpha) - 1$. Here we
argue exactly as in case (1) with $j = k(\alpha) - 1$.
\item Say $\beta \in I_m$ and $k(\beta) = k(\alpha)$. Here
the summand $N(U)_{k(\alpha)} \cdot \alpha$ is mapped by the
identity to the summand $N(U)_{k(\beta)} \cdot \beta$.
This is the same as the effect of $U(\varphi)$ since in
this case $U(\varphi) \circ U(\alpha) = U(\beta)$.
\item Say $\beta \in I_m$ and $k(\beta) = k(\alpha) - 1$.
Here we use the differential $(-1)^{k(\alpha)} d_{N(U), k(\alpha)}$
to map the summand $N(U)_{k(\alpha)} \cdot \alpha$
to the summand $N(U)_{k(\beta)} \cdot \beta$.
On the other hand, since $\Im(\beta) = [k(\beta)]$ in this case
we get $\alpha \circ \varphi = \delta^{k(\alpha)}_{k(\alpha)} \circ \beta$.
Thus we see that $U(\varphi)$ composed with the restriction of
$U(\alpha)$ to $N(U)_{k(\alpha)}$ is equal to
$U(\beta)$ precomposed with $d^{k(\alpha)}_{k(\alpha)}$ restricted to
$N(U)_{k(\alpha)}$.
Since $d_{N(U), k(\alpha)} = \sum (-1)^i d^{k(\alpha)}_i$
and since $d^{k(\alpha)}_i$ restricts to zero on
$N(U)_{k(\alpha)}$ for $i < k(\alpha)$
we see that equality holds.
\end{enumerate}
This finishes the proof of the theorem.
\end{proof}
